\input{../tex-inputs/latex-leseansicht-vorspann}

\section[Arthur Schnitzler und Hugo von Hofmannsthal an Gustav Schwarzkopf, 28. 5. 1899]{L04421 Arthur Schnitzler und Hugo von Hofmannsthal an Gustav Schwarzkopf, 28. 5. 1899}
\nopagebreak\mylabel{L04421v}
\rehead{ }\normalsize\beginnumbering\briefempfaengerindex{Schwarzkopf, Gustav@\textsc{Schwarzkopf, Gustav}!zzzHofmannsthal, Hugo von@\emph{von Hugo von Hofmannsthal}!1899-05-281@{28. 5. 1899}|(be}\briefempfaengerindex{Schwarzkopf, Gustav@\textsc{Schwarzkopf, Gustav}!zzzSchnitzler, Arthur@\emph{von Arthur Schnitzler}!1899-05-281@{28. 5. 1899}|(be}
\toendnotes[C]{\smallbreak\pagebreak[2]}
\correspDesc{Versand  durch Arthur Schnitzler, Hugo Hofmannsthal am 28. 5. 1899 in Rosenburg
\newline{}Übermittlung  am 28. 5. 1899 in Rosenburg
\newline{}Erhalt  durch Gustav Schwarzkopf im Zeitraum [29. 5. 1899 – 2. 6. 1899?] in Wien}\toendnotes[C]{\smallbreak}
\Standort{DLA, A:Schnitzler, HS.1985.1.1897.}
\physDesc{Bildpostkarte, 159 Zeichen
\newline{}Handschrift Arthur Schnitzler: Bleistift, deutsche Kurrent
\newline{}Handschrift Hugo von Hofmannsthal: Bleistift
\newline{}Versand: Stempel: »\nobreak{}\oindex{Bahnhof Rosenburg@\textbf{Bahnhof Rosenburg}|pwk}Rosenburg Bahnhof, 28/5 99\nobreak{}«.  }\pstart{}{\pb}Herrn \textsc{Gustav
                     Schwarzkopf}\pend{}\pstart{}Wien\oindex{Wien@\textbf{Wien}|pw}\pend{}\pstart{}\textsc{I. Tiefer Graben 23}\oindex{Wien@\textbf{Wien}!I., Innere Stadt@\textbf{I., Innere Stadt}!Tiefer Graben 23@\textbf{Tiefer Graben 23}|pw}\pend{}{\bigskip}
\pstart
           \centering{}{\pb}\textcolor{gray}{\textbf{Gruss aus Rosenburg im
                     Kampthal\oindex{Rosenburg@\textbf{Rosenburg}|pw}.}}\pend
           \vspace{1em}
\pstart
           \raggedleft{}{\pb}So{\geminationn}tag Abend\pend
           \vspace{0.5em}
\pstart
           \textsc{Sigmundsherberg\oindex{Sigmundsherberg@\textbf{Sigmundsherberg}|pw}} – – \textsc{Horn\oindex{Horn@\textbf{Horn}|pw}} – hieher geradelt – dieſes Bild ist geſchmeichelt.\pend
           \pstart Herzliche Grüße Ihr \spacefill\mbox{Arthur}\pend{}\selectlanguage{ngerman}\vspace{1em}\pstart \spacefill\mbox{{[}hs. Hofmannsthal:{]} Hugo}\pend{}\selectlanguage{ngerman}\endnumbering\briefempfaengerindex{Schwarzkopf, Gustav@\textsc{Schwarzkopf, Gustav}!zzzHofmannsthal, Hugo von@\emph{von Hugo von Hofmannsthal}!1899-05-281@{28. 5. 1899}|)be}\briefempfaengerindex{Schwarzkopf, Gustav@\textsc{Schwarzkopf, Gustav}!zzzSchnitzler, Arthur@\emph{von Arthur Schnitzler}!1899-05-281@{28. 5. 1899}|)be}\mylabel{L04421h}
\begin{anhang}
\end{anhang}\newcommand{\dateiname}{L04421}\newcommand{\titel}{Arthur Schnitzler und Hugo von Hofmannsthal an Gustav Schwarzkopf, 28. 5. 1899}\newcommand{\editorInnen}{Herausgegeben von Jahnke, SelmaMüller, Martin Anton}\input{../tex-inputs/latex-leseansicht-abspann}
